\lecture{The abstract framework: Lax--Milgram, C\'ea, and the move to 2D}{Lectures 13--18 --- late September to mid-October, 2026}{Math 6660 --- \texttt{prashant2.pdf} (colleague's notes)}

\noindent{\small\itshape These lectures are missing from my own notebook; the
transcription below follows a classmate's notes, so the wording and the order of
presentation are his, not mine. Only two dates are legible in the scan ---
\textbf{2/Oct/26} at the abstract variational formulation, and
\textbf{12/Oct/26} at the biharmonic example --- so the lecture numbering above
is my reconstruction, not something written in the source. Pages 1--18 of that
scan repeat material already transcribed in Lectures 7--12; the transcription
starts where the new material does.}

\begin{roadmap}[Lectures 13--18]
\roadq{Every problem so far has been solved by the same four moves. What is the
       general theorem, and what does it cost to go to two dimensions?}
\begin{roadlist}
\item The bar and the beam get rewritten in \emph{one} language:
  $B(u,v)=L(v)$. Everything specific to the problem is pushed into three
  properties of $B$ --- symmetric, bounded, elliptic.
\item \textbf{Lax--Milgram} then says: those three properties alone guarantee a
  unique solution. No more re-deriving existence for each new PDE.
\item \textbf{C\'ea's lemma} in general form: $\norm{u-u_h}\le
  (\beta^{*}/\beta_{*})\norm{u-v}$. The constant $1$ we had before was the lucky
  case.
\item The same abstract theorem gives, for free, that the stiffness matrix is
  \textbf{symmetric positive definite} --- in any dimension, any polynomial
  order, any mesh.
\item Then 2D: Poisson on a triangulated domain, the biharmonic plate (and why
  it needs the 21-degree-of-freedom Argyris element), and how a mesh is actually
  generated --- Voronoi, Delaunay, level sets.
\end{roadlist}
\payoff{This is the point where the course stops being about one-dimensional
examples and becomes the theory a real FE code is built on.}
\end{roadmap}

% ==================================================================
\section*{Part I --- the abstract formulation}
% ==================================================================
\subsection*{So far: two examples}

Both problems solved so far have the same shape.

\medskip\noindent
\textbf{(1) The two-point BVP.}
\begin{equation}\label{eq:l13ex1}
  -(au')'=f \ \text{ in } (0,1),\qquad u(0)=u(1)=0,
\end{equation}
\begin{equation}\label{eq:l13ex1V}
  \Vsp=\Bigl\{v:\ \int_0^1\bigl[v^{2}+(v')^{2}\bigr]\dx<\infty,\ \ v(0)=v(1)=0\Bigr\},
\end{equation}
\begin{equation}\label{eq:l13ex1B}
  B(u,v)=\int_0^1 a\,u'v'\dx,\qquad L(v)=\int_0^1 f v\dx .
\end{equation}

\noindent
\textbf{(2) The beam equation.}
\begin{equation}\label{eq:l13ex2V}
  \Vsp=\Bigl\{v:\ \int_0^1\bigl[v^{2}+(v')^{2}+(v'')^{2}\bigr]\dx<\infty,\ \
  v(0)=v(1)=v'(0)=v'(1)=0\Bigr\},
\end{equation}
\begin{equation}\label{eq:l13ex2B}
  B(u,v)=\int_0^1 u''v''\dx,\qquad L(v)=\int_0^1 f v\dx .
\end{equation}

\noindent
In both cases:
\begin{equation}\label{eq:l13absvar}
  \textbf{(Var.\ Prob.)}\quad\text{Find } u\in\Vsp \text{ s.t. }
  B(u,v)=L(v)\quad\forall\,v\in\Vsp ,
\end{equation}
\begin{equation}\label{eq:l13absmin}
  F(w)=\tfrac12 B(w,w)-L(w),\qquad
  \textbf{(Energy Min.)}\quad\text{Find } u\in\Vsp \text{ s.t. }
  F(u)\le F(w)\ \ \forall\,w\in\Vsp .
\end{equation}

\medskip\noindent
\textbf{Note.} We seek a \textbf{conforming} FEM: $\Vh\subset\Vsp$.
(A non-conforming FEM has $\Vh\not\subset\Vsp$.)

\begin{intuition}[what is actually being abstracted]
Nothing new is being claimed here --- \eqref{eq:l13ex1B} and
\eqref{eq:l13ex2B} are the weak forms already derived. The move is to stop
writing the integrals and give them a name. Once the bar and the beam are both
``$B(u,v)=L(v)$'', any theorem proved about $B$ and $L$ applies to both at
once, and to the 2D problems later in these lectures that have not been written
down yet.

The price is that $B$ and $L$ say nothing about \emph{which} problem you are
solving. Everything that distinguishes a bar from a beam has been pushed into
the space $\Vsp$ and into the three properties of $B$ defined below.
\end{intuition}

% ==================================================================
\subsection*{Multi-index notation}

Consider a domain $\Om\subset\R^{d}$. Derivatives are defined via a
\textbf{multi-index}: let $\alpha=(\alpha_{1},\alpha_{2},\dots,\alpha_{d})$ be a
$d$-tuple, where the $\alpha_{i}$ are non-negative integers. The \emph{length}
of $\alpha$ is
\begin{equation}\label{eq:l13mlen}
  |\alpha|=\alpha_{1}+\alpha_{2}+\dots+\alpha_{d},
\end{equation}
and the corresponding derivative of $v$ is
\begin{equation}\label{eq:l13Dalpha}
  D^{\alpha}v=\frac{\partial^{|\alpha|}v}
       {\partial x_{1}^{\alpha_{1}}\,\partial x_{2}^{\alpha_{2}}\cdots
        \partial x_{d}^{\alpha_{d}}} .
\end{equation}

\noindent
\textbf{Ex.} If $d=2$ and $|\alpha|=2$:
\[
  D^{(2,0)}v=\frac{\partial^{2}v}{\partial x_{1}^{2}},\qquad
  D^{(0,2)}v=\frac{\partial^{2}v}{\partial x_{2}^{2}},\qquad
  D^{(1,1)}v=\frac{\partial^{2}v}{\partial x_{1}\partial x_{2}} .
\]

\begin{intuition}[why bother]
In 1D you can write $v$, $v'$, $v''$ and be done. In 2D ``the second
derivatives'' is already three different objects, and in $d$ dimensions it is
$\binom{d+1}{2}$ of them. The multi-index is bookkeeping that lets
``\emph{all} derivatives up to order $r$'' be written as a single sum
$\sum_{|\alpha|\le r}$ instead of an expanding list. Every Sobolev definition
below uses exactly that.
\end{intuition}

% ==================================================================
\subsection*{Sobolev spaces}

\begin{equation}\label{eq:l13Lloc}
  L^{1}_{\mathrm{loc}}(\Om)\defeq\bigl\{v:\ v\in L^{1}(K)
  \text{ for any compact subset } K \text{ inside } \Om\bigr\} .
\end{equation}

\noindent
\textbf{Defn (Sobolev space).}
\begin{equation}\label{eq:l13Wrq}
  W^{r,q}(\Om)=\Bigl\{v\in L^{1}_{\mathrm{loc}}(\Om):\
  \norm{v}_{W^{r,q}(\Om)}<\infty\Bigr\},
\end{equation}
where
\begin{equation}\label{eq:l13Wnorm}
  \norm{v}_{W^{r,q}(\Om)}\defeq
  \Bigl[\sum_{|\alpha|\le r}\norm{D^{\alpha}v}^{q}_{L^{q}(\Om)}\Bigr]^{1/q},
  \qquad 1\le q<\infty,\quad r\ge0,
\end{equation}
and
\begin{equation}\label{eq:l13Lqnorm}
  \norm{v}^{q}_{L^{q}(\Om)}=\int_{\Om}|v|^{q}\dx .
\end{equation}

\noindent
\textbf{Defn (Sobolev space with homogeneous BCs).}
\begin{equation}\label{eq:l13W0}
  W^{r,q}_{0}(\Om)=\Bigl\{v\in W^{r,q}(\Om):\
  D^{\alpha}v\big|_{\dOm}=0 \text{ for } |\alpha|<r\Bigr\} .
\end{equation}

\todoq{On the page where the $L^{q}$ norm is first introduced my colleague's
notes write $\norm{u}_{L^{q}(\Om)}\defeq\int_{\Om}|u|\dx$, with no exponent $q$
and no outer root. That cannot be right --- it would make the $L^{q}$ norm
independent of $q$ --- and the correct form appears two lines later as
\eqref{eq:l13Lqnorm}. Transcribed as \eqref{eq:l13Lqnorm}; worth confirming
against the lecturer's own notes.}

% ==================================================================
\subsection*{Hilbert spaces}

\textbf{Defn (Hilbert space).} When $q=2$ we write
$W^{r,2}(\Om)=\Hs{r}$. This is called a \textbf{Hilbert space}, with
\begin{equation}\label{eq:l13Hnorm}
  \norm{v}_{\Hs{r}}=\sqrt{(v,v)_{\Hs{r}}},
  \qquad
  (v,w)_{\Hs{r}}\defeq\sum_{|\alpha|\le r}\int_{\Om}
     \bigl(D^{\alpha}v\bigr)\bigl(D^{\alpha}w\bigr)\dx .
\end{equation}

\medskip\noindent
\textbf{Note.} $r=0$ also makes sense: $H^{0}(\Om)=\Ltwo$.

\medskip\noindent
\textbf{Defn (semi-norm).}
\begin{equation}\label{eq:l13semi}
  |v|_{\Hs{r}}=\sqrt{\sum_{|\alpha|=r}\int_{\Om}\bigl(D^{\alpha}v\bigr)^{2}\dx} .
\end{equation}

\begin{intuition}[norm vs.\ semi-norm, in one line]
The norm sums over $|\alpha|\le r$ --- \emph{every} derivative up to order $r$,
including the function itself. The semi-norm sums over $|\alpha|=r$ --- only the
\emph{top} ones. So $|v|_{H^{r}}=0$ does not force $v=0$; it only forces $v$ to
be a polynomial of degree $<r$. That is why a semi-norm appears on the
right-hand side of every interpolation estimate: what the interpolant cannot
reproduce is exactly the top derivative, and the lower ones are irrelevant to
the error.
\end{intuition}

\begin{dictionary}[the spaces in engineering terms]
$\Ltwo$ & finite ``size'': $\int v^{2}<\infty$\\
$\Hone$ & finite strain energy for a bar: $\int (v')^{2}<\infty$\\
$\Honez$ & the same, and clamped at the boundary\\
$H^{2}_{0}(\Om)$ & finite bending energy, clamped \emph{and} zero slope --- the beam/plate space\\
$|v|_{H^{r}}$ & the energy-like part only, with the rigid-body content stripped out\\
\end{dictionary}

\medskip\noindent
\textbf{Note.} For the remainder of the discussion, let $\Vsp$ be a Hilbert
space with inner product $(\cdot,\cdot)_{\Vsp}$ and corresponding norm
$\norm{\cdot}_{\Vsp}$.

% ==================================================================
\subsection*{Linear functionals and the dual space}

\textbf{Defn.} A mapping $L:\Vsp\to\R$ is a \textbf{linear functional} if
\begin{equation}\label{eq:l13lin}
  L(\alpha u+\beta v)=\alpha L(u)+\beta L(v)
  \qquad\forall\,\alpha,\beta\in\R,\ \forall\,u,v\in\Vsp .
\end{equation}
$L$ is \textbf{bounded} in the norm $\norm{\cdot}_{\Vsp}$ if there exists a
constant $\widetilde{L}>0$ such that
\begin{equation}\label{eq:l13bdd}
  |L(v)|\le\widetilde{L}\,\norm{v}_{\Vsp}\qquad\forall\,v\in\Vsp .
\end{equation}

\medskip\noindent
\textbf{Note.} A bounded linear functional is continuous:
\[
  L(v)-L(w)=L(v-w)
  \ \Longrightarrow\
  |L(v)-L(w)|=|L(v-w)|\le\widetilde{L}\norm{v-w}_{\Vsp} .
\]
The converse is also true: a continuous linear functional is bounded.

\medskip\noindent
\textbf{Defn.} The set of bounded linear functionals on $\Vsp$ is termed the
\textbf{dual space} of $\Vsp$, denoted $\Vsp'$, equipped with the norm
\begin{equation}\label{eq:l13dualnorm}
  \norm{L}_{\Vsp'}=\sup_{\substack{v\in\Vsp\\ v\neq0}}
  \frac{L(v)}{\norm{v}_{\Vsp}}\;\le\;\widetilde{L} .
\end{equation}

\medskip\noindent
\textbf{Ex.\ (dual space of $L^{q}(\Om)$).} Consider the linear functional
$L(v)=\int_{\Om}f(x)v(x)\dx$ for $f\in L^{q'}(\Om)$ where
$\tfrac1q+\tfrac1{q'}=1$. \textbf{H\"older's inequality},
\begin{equation}\label{eq:l13holder}
  \norm{uv}_{L^{1}(\Om)}\le\norm{u}_{L^{q'}(\Om)}\norm{v}_{L^{q}(\Om)}
  \qquad\forall\,u\in L^{q'}(\Om),\ v\in L^{q}(\Om),
\end{equation}
gives $|L(v)|\le\norm{f}_{L^{q'}}\norm{v}_{L^{q}}$. So every
$f\in L^{q'}(\Om)$ can be viewed as a bounded linear functional on $L^{q}(\Om)$;
every bounded linear functional on $L^{q}(\Om)$ is of this form; hence
$L^{q'}(\Om)$ is the dual space of $L^{q}(\Om)$.

\medskip\noindent
\textbf{Thm (Riesz representation).} Let $H$ be a Hilbert space with scalar
product $(\cdot,\cdot)_{H}$. Then for any continuous linear functional
$L:H\to\R$ there exists a \emph{unique} $u\in H$ such that
\begin{equation}\label{eq:l13riesz}
  L(v)=(u,v)_{H}\qquad\forall\,v\in H .
\end{equation}

\begin{mechanics}[what Riesz is doing for us]
A load is a functional: it eats a virtual displacement and returns a number
(the work done). Riesz says that any such load, provided it is bounded, is
secretly \emph{a displacement field} paired against the virtual displacement
through the energy inner product. That identification is the entire existence
proof for Lax--Milgram below --- the load $L$ gets converted into the solution
$u$.
\end{mechanics}

% ==================================================================
\subsection*{Abstract variational problems}

\textbf{Defn (bilinear form).} Let $\Vsp$ be a Hilbert space with scalar
product $(\cdot,\cdot)_{\Vsp}$. Then $B(\cdot,\cdot):\Vsp\times\Vsp\to\R$ is a
\textbf{bilinear form} if
\begin{equation}\label{eq:l13bilin}
  B(u,\alpha v+\beta w)=\alpha B(u,v)+\beta B(u,w),\qquad
  B(\alpha u+\beta v,w)=\alpha B(u,w)+\beta B(v,w)
\end{equation}
for all $\alpha,\beta\in\R$ and all $u,v,w\in\Vsp$.

\medskip
The three properties that everything below depends on:

\medskip\noindent\fbox{\parbox{0.955\linewidth}{%
\textbf{(S) Symmetric.} $B(u,v)=B(v,u)\quad\forall\,u,v\in\Vsp$.

\medskip
\textbf{(CB) Continuous / bounded.} There exists $\beta^{*}>0$ such that
\[
  B(u,v)\le|B(u,v)|\le\beta^{*}\norm{u}_{\Vsp}\norm{v}_{\Vsp}
  \qquad\forall\,u,v\in\Vsp .
\]

\medskip
\textbf{(VE) $\Vsp$-elliptic (a.k.a.\ coercive).} There exists $\beta_{*}>0$
such that
\[
  B(v,v)\ge\beta_{*}\norm{v}^{2}_{\Vsp}\qquad\forall\,v\in\Vsp .
\]}}

\medskip\noindent
\textbf{Defn (abstract variational problem).} Let
$B(\cdot,\cdot):\Vsp\times\Vsp\to\R$ be a bilinear form satisfying (S), (CB)
and (VE). Let $L(\cdot):\Vsp\to\R$ be a linear functional bounded in
$\norm{\cdot}_{\Vsp}$, \ie $|L(v)|\le\widetilde{L}\norm{v}_{\Vsp}$
$\forall v\in\Vsp$. Then
\begin{align}
  \textbf{(1) Var.\ Prob.}&\quad\text{Find } u\in\Vsp \text{ s.t. }
     B(u,v)=L(v)\ \ \forall\,v\in\Vsp,\label{eq:l13VP}\\
  \textbf{(2) Energy Min.}&\quad\text{Find } u\in\Vsp \text{ s.t. }
     F(u)\le F(w)\ \ \forall\,w\in\Vsp,\quad
     F(v)=\tfrac12 B(v,v)-L(v).\label{eq:l13EM}
\end{align}

\begin{mechanics}[the three conditions as structural statements]
\textbf{(S)} is Maxwell--Betti reciprocity: the work done by load A through
displacement B equals the work of B through A. It is what makes $\Kmat$
symmetric.

\textbf{(CB)} says the structure cannot store infinite energy from a finite
displacement --- the material is bounded in stiffness.

\textbf{(VE)} says every admissible non-zero displacement stores
\emph{positive} energy, bounded below by $\beta_{*}$. This is the statement
that the structure has no mechanism and no rigid-body mode left: it is why
$\Kmat$ is positive definite and not merely positive semi-definite.
$\beta_{*}$ is, in effect, the softest mode of the structure.
\end{mechanics}

% ------------------------------------------------------------------
\subsubsection*{Claim: (1) and (2) are equivalent}

\paragraph{($\Longrightarrow$)} Assume $u\in\Vsp$ satisfies the variational
problem. Let $w=u+v\in\Vsp$ for some $v\in\Vsp$. Then
\begin{align}
  F(w)=F(u+v)&=\tfrac12 B(u+v,u+v)-L(u+v),\notag\\
  B(u+v,u+v)&=B(u,u+v)+B(v,u+v)=B(u,u)+2B(u,v)+B(v,v)
     \quad\text{(used (S))},\label{eq:l13expand}
\end{align}
so
\begin{equation}\label{eq:l13complete}
  F(u+v)=\underbrace{\Bigl[\tfrac12 B(u,u)-L(u)\Bigr]}_{F(u)}
   +\underbrace{\Bigl[B(u,v)-L(v)\Bigr]}_{=\,0\ \text{by (Var.\ Prob.)}}
   +\tfrac12 B(v,v),
\end{equation}
\ie $F(w)=F(u)+\tfrac12 B(v,v)$. But $B(v,v)\ge\beta_{*}\norm{v}^{2}_{\Vsp}>0$
for $v\neq0$ by (VE), hence
\begin{equation}
  F(u)\le F(w)\qquad\forall\,w\in\Vsp .
\end{equation}

\paragraph{($\Longleftarrow$)} Assume $u\in\Vsp$ satisfies $F(u)\le F(w)$
$\forall w\in\Vsp$. Let
\begin{equation}\label{eq:l13g}
  g(\eps)=F(u+\eps v)\qquad\text{for fixed but arbitrary } v\in\Vsp .
\end{equation}
Expanding,
\begin{equation}\label{eq:l13gexp}
  g(\eps)=F(u)+\eps B(u,v)-\eps L(v)+\tfrac{\eps^{2}}{2}B(v,v),
\end{equation}
\begin{equation}\label{eq:l13gprime}
  g'(\eps)=B(u,v)-L(v)+\eps\,B(v,v),
\end{equation}
and $g'(0)=0$ gives $B(u,v)=L(v)$. \qed

\begin{recall}[this is Lecture 3 and Lecture 11 again]
The bar version of this proof is Lecture 3; the beam version is Lecture 11.
The algebra is identical in all three --- complete the square one way, take
$\tfrac{\dd}{\dd\eps}$ at $\eps=0$ the other. What is new is only that it is
now done \emph{once}, for every problem of this shape at the same time.
\end{recall}

% ==================================================================
\subsection*{The Lax--Milgram theorem \hfill {\normalfont\small [2/Oct/26]}}

\medskip\noindent\fbox{\parbox{0.955\linewidth}{%
\textbf{Thm (Lax--Milgram).} Let $\Vsp$ be a Hilbert space with norm
$\norm{\cdot}_{\Vsp}$. Let $B(\cdot,\cdot):\Vsp\times\Vsp\to\R$ be a bilinear
form satisfying (S), (CB) and (VE). Let $L(\cdot):\Vsp\to\R$ be a bounded
linear functional. Then there exists a \textbf{unique} solution to the
variational problem
\[
  \text{find } u\in\Vsp \text{ s.t. } B(u,v)=L(v)\ \ \forall\,v\in\Vsp,
\]
with the property that
\[
  \norm{u}_{\Vsp}\le\frac{\widetilde{L}}{\beta_{*}},
\]
where $\widetilde{L},\beta_{*}$ are real positive constants.}}

\medskip\noindent
\textbf{Proof.} Since $B(\cdot,\cdot)$ is symmetric and $\Vsp$-elliptic, it
induces a scalar product on $\Vsp$:
\begin{equation}\label{eq:l13energyip}
  [u,v]\defeq B(u,v)\qquad\forall\,u,v\in\Vsp .
\end{equation}
By (VE) and (CB),
\begin{equation}\label{eq:l13equiv}
  [v,v]=B(v,v)\ge\beta_{*}\norm{v}^{2}_{\Vsp},
  \qquad
  [v,v]=B(v,v)\le\beta^{*}\norm{v}^{2}_{\Vsp},
\end{equation}
so
\begin{equation}\label{eq:l13equiv2}
  \boxed{\;\beta_{*}\norm{v}^{2}_{\Vsp}\le[v,v]\le\beta^{*}\norm{v}^{2}_{\Vsp}\;}
\end{equation}
and the norms $\sqrt{[v,v]}$ and $\norm{v}_{\Vsp}$ are \textbf{equivalent}.
Therefore $L(\cdot)$ is bounded in both. By the Riesz representation theorem
\eqref{eq:l13riesz} there exists a unique $u\in\Vsp$ such that
$L(v)=[u,v]$ $\forall v\in\Vsp$, \ie a unique $u\in\Vsp$ with
$B(u,v)=L(v)$ $\forall v\in\Vsp$.

\medskip
For the solution bound,
\begin{equation}\label{eq:l13bound}
  \beta_{*}\norm{u}^{2}_{\Vsp}\le B(u,u)=L(u)\le\widetilde{L}\norm{u}_{\Vsp}
  \ \Longrightarrow\
  \norm{u}_{\Vsp}\le\frac{\widetilde{L}}{\beta_{*}} . \qed
\end{equation}

\begin{intuition}[what Lax--Milgram buys, in one sentence]
Up to now, every new BVP needed its own argument that a solution exists and is
unique. Lax--Milgram replaces all of them with a checklist: write the weak
form, verify (S), (CB), (VE), and you are done --- existence, uniqueness, and a
bound on how big the answer can get, all three.

The bound \eqref{eq:l13bound} is also the \emph{stability} statement: double
the load and the response at most doubles, with constant $1/\beta_{*}$. A
structure with a nearly-singular mode has a tiny $\beta_{*}$ and is predicted,
correctly, to be sensitive.
\end{intuition}

\begin{pitfall}[(S) is not actually needed for Lax--Milgram]
The theorem as stated in lecture assumes (S), (CB) \emph{and} (VE), and the
proof given uses (S) to build the inner product $[u,v]$ and then quotes Riesz.
The standard Lax--Milgram theorem does \emph{not} require symmetry --- it holds
for non-symmetric $B$, at the cost of a longer proof that cannot use Riesz
directly. Symmetry is what lets the short proof above work, and it is also what
gives the energy-minimisation reading \eqref{eq:l13EM}; a non-symmetric problem
(\eg\ one with a $b\,u'$ convection term, cf.\ Lecture 1) has no energy
functional to minimise but still has a unique weak solution.
\end{pitfall}

% ==================================================================
\section*{Part II --- conforming FEM in the abstract setting}
% ==================================================================
\subsection*{The discrete problem}

Replace $\Vsp$ by a finite-dimensional subspace $\Vh\subset\Vsp$
(\emph{conforming}):
\begin{equation}\label{eq:l13fem}
  \text{Find } \uh\in\Vh \text{ s.t. } B(\uh,v)=L(v)\quad\forall\,v\in\Vh .
\end{equation}
Lax--Milgram still applies to this, so there exists a unique solution with
$\norm{\uh}_{\Vsp}\le\widetilde{L}/\beta_{*}$.

\medskip
Let $\{\phih_{i}\}_{i=1}^{M}$ be a basis for $\Vh$, take $v=\phih_{j}$,
$j=1,\dots,M$, and write $\uh=\sum_{i=1}^{M}\xi_{i}\phih_{i}$. Then
\begin{equation}\label{eq:l13sum}
  \sum_{i=1}^{M}B(\phih_{i},\phih_{j})\,\xi_{i}=L(\phih_{j}),
  \qquad j=1,2,\dots,M,
\end{equation}
\ie the linear system
\begin{equation}\label{eq:l13system}
  \boxed{\;\bm{A}\,\bm{\xi}=\bm{b}\;}
  \qquad
  A_{ij}=B(\phih_{i},\phih_{j}),\quad
  \bm{\xi}=(\xi_{1},\xi_{2},\dots,\xi_{M}),\quad
  b_{j}=L(\phih_{j}).
\end{equation}

% ------------------------------------------------------------------
\subsubsection*{Claim: under (S) and (VE), $\bm{A}$ is symmetric positive definite}

\paragraph{(1) Symmetry} follows directly from (S):
\begin{equation}
  A_{ij}=B(\phih_{i},\phih_{j})=B(\phih_{j},\phih_{i})=A_{ji}.
\end{equation}

\paragraph{(2) Positive definiteness.} Let
$w=\sum_{i=1}^{M}w_{i}\phih_{i}\in\Vh$. By (VE),
$B(w,w)\ge\beta_{*}\norm{w}^{2}_{\Vsp}$, so if $w\not\equiv0$ then
$B(w,w)>0$. Expanding,
\begin{equation}\label{eq:l13quad}
  B\Bigl(\sum_{i=1}^{M}w_{i}\phih_{i},\,w\Bigr)
   =\sum_{i=1}^{M}w_{i}B(\phih_{i},w)
   =\sum_{i=1}^{M}w_{i}B\Bigl(\phih_{i},\sum_{j=1}^{M}w_{j}\phih_{j}\Bigr)
   =\sum_{i=1}^{M}\sum_{j=1}^{M}w_{i}\underbrace{B(\phih_{i},\phih_{j})}_{A_{ij}}w_{j}
   =\bm{w}^{T}\bm{A}\,\bm{w}>0
\end{equation}
by (VE). Hence $\bm{A}$ is positive definite. \qed

\medskip\noindent
\textbf{Note.} This means a conforming FEM for this abstract variational
problem will generate an SPD linear system --- \emph{independent of} the number
of dimensions, the polynomial order, the shape of the solution domain $\Om$,
the type of mesh, etc.

\begin{intuition}[why this note is the point of the whole abstraction]
Every stiffness matrix met so far came out symmetric positive definite, and
each time it was verified by looking at the numbers. \eqref{eq:l13quad} shows
it was never a coincidence: $\bm{w}^{T}\bm{A}\bm{w}$ \emph{is} $B(w,w)$, the
energy of the function $w$, so (VE) --- a statement about the physics ---
becomes positive definiteness, a statement about the matrix. Nothing about 1D,
cubics or uniform meshes entered the argument, which is exactly why the result
survives into 2D in Part III.
\end{intuition}

% ==================================================================
\subsection*{C\'ea's lemma}

\medskip\noindent\fbox{\parbox{0.955\linewidth}{%
\textbf{Thm (C\'ea's lemma).} Let $\Vsp$ be a Hilbert space with norm
$\norm{\cdot}_{\Vsp}$, let $B$ be a bilinear form satisfying (S), (CB) and
(VE), and let $L$ be a linear bounded functional. Let $\Vh\subset\Vsp$ be a
finite-dimensional subspace, and let $u$, $\uh$ solve
\eqref{eq:l13VP} and \eqref{eq:l13fem} respectively. Then
\[
  \norm{u-\uh}_{\Vsp}\;\le\;\frac{\beta^{*}}{\beta_{*}}\,\norm{u-v}_{\Vsp}
  \qquad\forall\,v\in\Vh .
\]}}

\medskip\noindent
\textbf{Proof.} Subtract (Var.\ Prob.) and (FEM):
\begin{equation}\label{eq:l13GO}
  B(u,v)-B(\uh,v)=0\ \ \forall\,v\in\Vh
  \qquad\Longrightarrow\qquad
  \boxed{\;B(u-\uh,\,v)=0\quad\forall\,v\in\Vh\;}
\end{equation}
which is \textbf{Galerkin orthogonality}. Now take
\begin{equation}\label{eq:l13cea1}
  B(u-\uh,\,u-\uh)=B(u-\uh,\,u-v+v-\uh)
   =B(u-\uh,\,u-v)+\underbrace{B(u-\uh,\,v-\uh)}_{=\,0\ \text{by \eqref{eq:l13GO}}} ,
\end{equation}
so, using (VE) on the left and (CB) on the right,
\begin{equation}\label{eq:l13cea2}
  \beta_{*}\norm{u-\uh}^{2}_{\Vsp}\le B(u-\uh,u-\uh)=B(u-\uh,u-v)
   \le\beta^{*}\norm{u-\uh}_{\Vsp}\norm{u-v}_{\Vsp},
\end{equation}
and cancelling one factor $\norm{u-\uh}_{\Vsp}$,
\begin{equation}\label{eq:l13cea3}
  \boxed{\;\norm{u-\uh}_{\Vsp}\le\frac{\beta^{*}}{\beta_{*}}\norm{u-v}_{\Vsp}\;}
  \qquad\qed
\end{equation}

\medskip\noindent
\textbf{Note (error estimate in $\norm{\cdot}_{\Vsp}$).} Take
$v=\Ih u$, the projection of $u\in\Vsp$ into $\Vh$. Then
\begin{equation}\label{eq:l13ceainterp}
  \norm{u-\uh}_{\Vsp}\le\frac{\beta^{*}}{\beta_{*}}\norm{u-\Ih u}_{\Vsp},
\end{equation}
where $\norm{u-\Ih u}_{\Vsp}$ is the \emph{interpolation} error measured in
$\norm{\cdot}_{\Vsp}$. If you can generate a sequence of meshes and associated
bases such that $\norm{u-\Ih u}_{\Vsp}\to0$ as $h\to0$, then you have
convergence; typically
\begin{equation}\label{eq:l13rate}
  \norm{u-\Ih u}_{\Vsp}=\mathcal{O}(h^{p}).
\end{equation}

\begin{pitfall}[the constant is no longer $1$]
Lecture 6 proved $\norm{u-U}_{E}\le\norm{u-v}_{E}$ --- constant exactly $1$ ---
and that is \emph{not} the general case. It happened because the energy norm
there was built from the same $a$ as the bilinear form, making
$\beta^{*}=\beta_{*}=1$ by construction. In a general norm
$\norm{\cdot}_{\Vsp}$ the constant is $\beta^{*}/\beta_{*}$, a condition-number
sort of quantity: $\uh$ is \emph{quasi}-optimal, within a fixed factor of the
best approximation, not the best approximation itself.
\end{pitfall}

% ==================================================================
\subsection*{Example: the 1D Dirichlet problem, checked against the three conditions}

\begin{equation}\label{eq:l13ex1d}
  -u''=f \ \text{ in } (0,1),\qquad u(0)=u(1)=0,
  \qquad I=(0,1),\qquad \Vsp=H^{1}_{0}(I),
\end{equation}
\begin{equation}\label{eq:l13ex1dnorm}
  \norm{v}^{2}_{\Vsp}=\int_{0}^{1}\bigl[(v')^{2}+v^{2}\bigr]\dx
   =\norm{v'}^{2}_{L^{2}(I)}+\norm{v}^{2}_{L^{2}(I)},
\end{equation}
\begin{equation}\label{eq:l13ex1dBL}
  B(u,v)=\int_{0}^{1}u'v'\dx,\qquad L(v)=\int_{0}^{1}fv\dx .
\end{equation}

\noindent
\textbf{Check conditions.}
\begin{enumerate}[label=\arabic*),leftmargin=2.2em,itemsep=3pt]
\item \textbf{Symmetric:} $(u',v')=(v',u')$. \cmark
\item \textbf{Continuous / bounded:}
  $|B(u,v)|\le\norm{u'}_{L^{2}(I)}\norm{v'}_{L^{2}(I)}
   \le\norm{u}_{\Vsp}\norm{v}_{\Vsp}$, so $\beta^{*}=1$. \cmark
\item \textbf{$\Vsp$-elliptic:}
  $B(v,v)=\norm{v'}^{2}_{L^{2}(I)}
   \ge\tfrac12\bigl(\norm{v'}^{2}_{L^{2}(I)}+\norm{v}^{2}_{L^{2}(I)}\bigr)
   =\tfrac12\norm{v}^{2}_{\Vsp}$, using \textbf{Poincar\'e}
  $\norm{v}_{L^{2}(I)}\le\norm{v'}_{L^{2}(I)}$ for $v\in H^{1}_{0}(I)$; so
  $\beta_{*}=\tfrac12$. \cmark
\item \textbf{$L$ continuous:}
  $|L(v)|\le\norm{f}_{L^{2}(I)}\norm{v}_{L^{2}(I)}
   \le\norm{f}_{L^{2}(I)}\norm{v}_{\Vsp}$, so
  $\widetilde{L}=\norm{f}_{L^{2}(I)}$. \cmark
\end{enumerate}

\noindent
$\Rightarrow$ existence, uniqueness and boundedness, and C\'ea's lemma gives
$\norm{u-\uh}_{\Vsp}\le2\norm{u-v}_{\Vsp}$ for all $v\in\Vh$, hence
\begin{equation}\label{eq:l13cgrates}
  \text{CG}(1):\ \norm{u-\uh}_{\Vsp}\le c_{1}\norm{h\,u''}_{L^{2}(I)},
  \qquad
  \text{CG}(2):\ \norm{u-\uh}_{\Vsp}\le c_{2}\norm{h^{2}u'''}_{L^{2}(I)} .
\end{equation}

\begin{intuition}[Poincar\'e is what makes (VE) true]
(VE) asks $B(v,v)\ge\beta_{*}\norm{v}^{2}_{\Vsp}$, but $B(v,v)$ only sees
$v'$, while $\norm{v}^{2}_{\Vsp}$ also contains $v$ itself. Those are only
comparable because of Poincar\'e, which says a function pinned to zero at the
boundary cannot be large without having a large derivative. Remove the boundary
condition --- a free--free bar --- and Poincar\'e fails for constants,
$\beta_{*}=0$, and $\bm{A}$ becomes singular. The rigid-body mode shows up as
the failure of exactly this one inequality.
\end{intuition}

% ==================================================================
\section*{Part III --- two dimensions}
% ==================================================================
\subsection*{The 2D Poisson equation}

\begin{equation}\label{eq:l132dpoisson}
  -\lap u=-u_{,xx}-u_{,yy}=f \ \text{ on } \Om\subset\R^{2},
  \qquad u=0 \text{ on } \dOm .
\end{equation}

Multiply by a test function $v$ and integrate over $\Om$:
\begin{equation}\label{eq:l132dmult}
  -\iint_{\Om}v\,\lap u\dx\dd y=\iint_{\Om}f v\dx\dd y .
\end{equation}

\noindent
\textbf{Note (product rule).}
$\grad\!\cdot\!(v\grad u)=\grad v\cdot\grad u+v\,\lap u$, so
\begin{equation}\label{eq:l132dprod}
  -\iint_{\Om}v\,\lap u\dx\dd y
   =\iint_{\Om}\grad u\cdot\grad v\dx\dd y
   -\iint_{\Om}\grad\!\cdot\!(v\grad u)\dx\dd y .
\end{equation}
By the \textbf{divergence theorem},
\begin{equation}\label{eq:l13divthm}
  \iint_{\Om}\grad\!\cdot\!(v\grad u)\dx\dd y
   =\oint_{\dOm}v\,\grad u\cdot\nn\,\dd s,
  \qquad \nn=\text{outward unit normal to } \dOm .
\end{equation}
Assuming $v=0$ on $\dOm$, the boundary integral vanishes and
\begin{equation}\label{eq:l132dweak}
  \iint_{\Om}\grad u\cdot\grad v\dx\dd y=\iint_{\Om}f v\dx\dd y .
\end{equation}

\noindent
So
\begin{equation}\label{eq:l132dV}
  \Vsp=\Honez=\Bigl\{v:\ \iint_{\Om}\bigl[v^{2}+\|\grad v\|^{2}\bigr]\dx\dd y<\infty,
  \ \ v=0 \text{ on } \dOm\Bigr\},
  \qquad \|\grad v\|^{2}=(v_{,x})^{2}+(v_{,y})^{2},
\end{equation}
\begin{equation}\label{eq:l132dBL}
  B(u,v)\defeq\iint_{\Om}\grad u\cdot\grad v\dx\dd y,
  \qquad
  L(v)\defeq\iint_{\Om}f v\dx\dd y .
\end{equation}

\noindent
\textbf{Check conditions.}
\begin{enumerate}[label=\arabic*),leftmargin=2.2em,itemsep=3pt]
\item \textbf{Symmetric:} $B(u,v)=B(v,u)$. \cmark
\item \textbf{Continuous / bounded:}
  $|B(u,v)|\le\norm{\grad u}_{\Ltwo}\norm{\grad v}_{\Ltwo}
   \le\norm{u}_{\Vsp}\norm{v}_{\Vsp}$, so $\beta^{*}=1$. \cmark
\item \textbf{$\Vsp$-elliptic:}
  $B(v,v)=\norm{\grad v}^{2}_{\Ltwo}
   \ge\tfrac12\bigl(\norm{\grad v}^{2}_{\Ltwo}+\norm{v}^{2}_{\Ltwo}\bigr)$,
  using the \textbf{multi-dimensional Poincar\'e inequality}
  \begin{equation}\label{eq:l13poincare2d}
    \norm{v}_{\Ltwo}\le C\,|v|_{\Hs{1}}=C\norm{\grad v}_{\Ltwo}
    \qquad\forall\,v\in\Honez ,
  \end{equation}
  giving
  $B(v,v)\ge\tfrac12\min\bigl(\tfrac1{C^{2}},1\bigr)\norm{v}^{2}_{\Vsp}$,
  \ie $\beta_{*}=\tfrac12\min\bigl(\tfrac1{C^{2}},1\bigr)$, a constant
  depending only on $\dOm$. \cmark
\item \textbf{$L$ continuous:}
  $|L(v)|\le\norm{f}_{\Ltwo}\norm{v}_{\Ltwo}\le\norm{f}_{\Ltwo}\norm{v}_{\Vsp}$,
  so $\widetilde{L}=\norm{f}_{\Ltwo}$. \cmark
\end{enumerate}

\begin{intuition}[nothing in Part II needed changing]
The 2D problem was checked against the \emph{same} four lines as the 1D
problem, and passed. That is the return on the abstraction: Lax--Milgram gives
existence and uniqueness, the SPD claim gives the matrix structure, and C\'ea
gives the error bound --- none of them re-proved. The only genuinely
two-dimensional work was the divergence theorem in \eqref{eq:l13divthm} and
finding the right Poincar\'e constant.
\end{intuition}

% ==================================================================
\subsection*{CG(1) in 2D: the triangulation}

\textbf{Mesh $\Th$:} subdivide $\Om$ into a finite number of triangles
$K_{j}$, $j=1,2,\dots,M$, with
\begin{align}
  h_{j}&\defeq\text{longest edge of } K_{j}, \label{eq:l13hj}\\
  h&\defeq\max_{1\le j\le M}h_{j}, \label{eq:l13h}\\
  \rho_{j}&\defeq\text{diameter of the largest circle inscribed in } K_{j}.
  \label{eq:l13rhoj}
\end{align}

\noindent
\textbf{Defn.} $\Th$ is a \textbf{regular} triangulation if there exists a
constant $\alpha>0$, \emph{independent of $h$}, such that
\begin{equation}\label{eq:l13regular}
  \frac{h_{j}}{\rho_{j}}\le\alpha\qquad\forall\,j .
\end{equation}

\noindent
The discrete space is
\begin{equation}\label{eq:l132dVh}
  \Vh=\Bigl\{v\in\Hone:\ v\big|_{K_{j}}=a_{j}+b_{j}x+c_{j}y,\ \
  v\big|_{\dOm}=0\Bigr\},
\end{equation}
and C\'ea's lemma gives the error estimate
\begin{equation}\label{eq:l132derr}
  \norm{u-\uh}_{\Hs{1}}\le c\,\norm{h\,u}_{\Hs{2}},
  \qquad
  |v|^{2}_{\Hs{2}}=\iint_{\Om}\bigl[(v_{,xx})^{2}+2(v_{,xy})^{2}+(v_{,yy})^{2}\bigr]\dx\dd y .
\end{equation}

\begin{mechanics}[what regularity \eqref{eq:l13regular} is forbidding]
$h_{j}/\rho_{j}$ is an aspect ratio: longest edge over inscribed-circle
diameter. Bounding it uniformly forbids \emph{sliver} elements --- long thin
triangles that collapse towards a line as the mesh is refined. This is the
same element-quality check a commercial mesher reports, and the reason it
matters is that the interpolation estimate behind \eqref{eq:l132derr} carries
$1/\rho_{j}$: let the triangles go slivery and the constant $c$ blows up even
though $h$ is going to zero.
\end{mechanics}

% ==================================================================
\subsection*{The beam, in the abstract language}

The 1D fourth-order problem written as an abstract variational problem:
\begin{equation}\label{eq:l13beamabs}
  u^{(iv)}=f \ \text{ in } (0,1),\qquad u(0)=u(1)=u'(0)=u'(1)=0,
  \qquad \Vsp=H^{2}_{0}(I),\ \ I=(0,1),
\end{equation}
\begin{equation}\label{eq:l13beamnorm}
  \norm{v}_{\Vsp}=\norm{v}_{\Hs{2}[I]}
   =\Bigl(\int_{0}^{1}\bigl[v^{2}+(v')^{2}+(v'')^{2}\bigr]\dx\Bigr)^{1/2},
\end{equation}
\begin{equation}\label{eq:l13beamBL}
  B(u,v)=(u'',v''),\qquad L(v)=(f,v).
\end{equation}
Here $H^{2}_{0}(I)$ means $u=0$ \emph{and} $u'=0$ on the boundary.

\noindent
\textbf{Check conditions.}
\begin{enumerate}[label=\arabic*),leftmargin=2.2em,itemsep=3pt]
\item \textbf{Symmetric:} $(u'',v'')=(v'',u'')$. \cmark
\item \textbf{Continuous / bounded:}
  $|B(u,v)|\le\norm{u''}_{L^{2}(I)}\norm{v''}_{L^{2}(I)}
   \le\norm{u}_{\Vsp}\norm{v}_{\Vsp}$, $\beta^{*}=1$. \cmark
\item \textbf{$\Vsp$-elliptic:} $B(v,v)=\norm{v''}^{2}_{L^{2}(I)}$, and
  applying Poincar\'e twice,
  $\norm{v}_{L^{2}(I)}\le\norm{v'}_{L^{2}(I)}\le\norm{v''}_{L^{2}(I)}$
  for $v\in H^{2}_{0}(I)$, gives
  $B(v,v)\ge\tfrac13\bigl(\norm{v}^{2}_{L^{2}}+\norm{v'}^{2}_{L^{2}}
   +\norm{v''}^{2}_{L^{2}}\bigr)=\tfrac13\norm{v}^{2}_{\Vsp}$,
  so $\beta_{*}=\tfrac13$. \cmark
\item \textbf{$L$ continuous:} $\widetilde{L}=\norm{f}_{L^{2}(0,1)}$. \cmark
\end{enumerate}

\noindent
\textbf{CG(3) estimate} (C\'ea's lemma):
\begin{equation}\label{eq:l13cg3}
  \norm{u-\uh}_{\Hs{2}[(0,1)]}\le c\,\norm{h^{2}u}_{\Hs{4}[(0,1)]},
  \qquad
  |v|^{2}_{\Hs{4}[(0,1)]}=\int_{0}^{1}\bigl(v''''\bigr)^{2}\dx .
\end{equation}

\begin{recall}[this is HW3 Problems 4--5, done abstractly]
The homework proved $\norm{u-U}_{E}=\mathcal{O}(h^{2})$ by hand, with Galerkin
orthogonality, Cauchy and the interpolation estimate. \eqref{eq:l13cg3} is the
same result read straight off C\'ea's lemma with
$\beta^{*}/\beta_{*}=3$. The hand proof gave a better constant because it used
the energy norm, in which $\beta^{*}=\beta_{*}$; the abstract route gives the
same \emph{rate} with less work.
\end{recall}

% ==================================================================
\subsection*{The biharmonic equation \hfill {\normalfont\small [12/Oct/26]}}

The 2D plate --- the beam's two-dimensional counterpart:
\begin{equation}\label{eq:l13biharm}
  \lap^{2}u=\lap(\lap u)=f
  \ \Longrightarrow\
  u_{,xxxx}+2u_{,xxyy}+u_{,yyyy}=f ,
\end{equation}
\begin{equation}\label{eq:l13biharmbc}
  \text{BCs:}\qquad u\big|_{\dOm}=0,\qquad \grad u\cdot\nn\big|_{\dOm}=0 .
\end{equation}

Multiply by $v$ and integrate, letting $z=\lap u$:
\begin{align}
  \iint_{\Om}v\,\lap(\lap u)\dx\dd y&=\iint_{\Om}vf\dx\dd y,\notag\\
  \iint_{\Om}v\,\lap z\dx\dd y
   &=-\iint_{\Om}\grad v\cdot\grad z\dx\dd y
     +\oint_{\dOm}v\,\grad z\cdot\nn\,\dd s ,\label{eq:l13biharmibp}
\end{align}
and taking $v=0$, $\grad v\cdot\nn=0$ on $\dOm$,
\begin{equation}\label{eq:l13biharmweak}
  B(u,v)=\iint_{\Om}\lap u\,\lap v\dx\dd y,
  \qquad
  L(v)=\iint_{\Om}f v\dx\dd y,
  \qquad
  \Vsp=H^{2}_{0}(\Om).
\end{equation}

\medskip\noindent
\textbf{Note.} If $v\in H^{2}_{0}(\Om)$ then
\begin{equation}\label{eq:l13H2norm}
  \norm{v}^{2}_{\Hs{2}}=\norm{v}^{2}_{\Ltwo}+\norm{\grad v}^{2}_{\Ltwo}
   +\norm{v_{,xx}}^{2}_{\Ltwo}+2\norm{v_{,xy}}^{2}_{\Ltwo}
   +\norm{v_{,yy}}^{2}_{\Ltwo},
\end{equation}
and $v\in H^{2}_{0}(\Om)$ means $v=0$, $v_{,x}=0$, $v_{,y}=0$ on $\dOm$.

\noindent
\textbf{Check conditions.}
\begin{enumerate}[label=\arabic*),leftmargin=2.2em,itemsep=3pt]
\item \textbf{Symmetric:} $B(u,v)=B(v,u)$. \cmark
\item \textbf{Continuous / bounded:}
  $|B(u,v)|\le\norm{\lap u}_{\Ltwo}\norm{\lap v}_{\Ltwo}
   \le\norm{u}_{\Vsp}\norm{v}_{\Vsp}$. \cmark
\item \textbf{$\Vsp$-elliptic:} $B(v,v)=\norm{\lap v}^{2}_{\Ltwo}$, and by a
  result of \textbf{Girault--Raviart (1981)},
  \begin{equation}\label{eq:l13girault}
    \norm{v}_{\Hs{2}}\le C\,\norm{\lap v}_{\Ltwo}
    \qquad\text{for } v\in H^{2}_{0}(\Om),
  \end{equation}
  hence $B(v,v)\ge\tfrac1{C^{2}}\norm{v}^{2}_{\Vsp}$. \cmark
\item \textbf{$L$ continuous:} $\widetilde{L}=\norm{f}_{\Ltwo}$. \cmark
\end{enumerate}

% ------------------------------------------------------------------
\subsubsection*{Conforming FEM for the biharmonic: the Argyris element}

A conforming method needs $\Vh\subset\Vsp=H^{2}_{0}(\Om)$, \ie
\begin{equation}\label{eq:l13argyrisV}
  \Vh=\Bigl\{v\in H^{2}(\Om):\
  v\big|_{K_{j}}=\sum_{\ell+m=0}^{\ell+m\,\le\,5}a^{(\ell m)}_{j}\,x^{(\ell)}y^{(m)},
  \ \ v\big|_{\dOm}=0,\ \ \grad v\cdot\nn\big|_{\dOm}=0\Bigr\},
\end{equation}
a complete quintic in two variables --- \textbf{21 different polynomials} on
each triangle. The corresponding element is the \textbf{Argyris element
(1968)}, whose $21$ degrees of freedom per triangle are
\begin{equation}\label{eq:l13argyrisdofs}
  \underbrace{u,\ \grad u,\ u_{,xx},\ u_{,xy},\ u_{,yy}}_{6\ \text{per vertex}\ \times\ 3\ \text{vertices}\ =\ 18}
  \qquad\text{plus}\qquad
  \underbrace{\grad u\cdot\nn}_{1\ \text{per edge}\ \times\ 3\ \text{edges}\ =\ 3}.
\end{equation}

\begin{mechanics}[why plates are hard, in one element]
A conforming plate element must be $C^{1}$ across element boundaries --- the
2D version of the beam's requirement --- and matching \emph{both} the value and
the full gradient along a shared edge is far more demanding than matching them
at a point. The cheapest complete polynomial that can do it is a quintic, which
is why the Argyris element carries second derivatives as nodal unknowns and
ends up with $21$ DOFs on a single triangle.

This is also why real plate and shell codes usually do \emph{not} use
conforming elements: they drop to non-conforming or mixed formulations (and
give up $\Vh\subset\Vsp$) rather than pay this price.
\end{mechanics}

% ==================================================================
\section*{Part IV --- 2D mesh generation}
% ==================================================================
\subsection*{Triangulations}

\textbf{Note.} FEM has no inherent restriction on the mesh. The simplest and
most flexible choice is to use a \textbf{simplex}: 2D triangles, 3D
tetrahedra.

\medskip\noindent
\textbf{Note.} A general triangulation in 2D is \textbf{unstructured}, meaning
grid points do not lie along coordinate lines.

\medskip\noindent
\textbf{Defn (triangulation).} Let $\Om$ be a 2D domain with polygonal
boundary $\dOm$. A triangulation $\Th=\{K_{j}\}_{j=1}^{M}$ is a subdivision of
$\Om$ into $M$ non-overlapping triangles constructed so that no vertex of one
triangle lies on the edge of another triangle:
\begin{equation}\label{eq:l13tri}
  \bigcup_{j=1}^{M}K_{j}=\Om .
\end{equation}
Let $X_{h}=\{x_{i}=(x_{i},y_{i})\}_{i=1}^{N}$ denote the set of \textbf{nodes}
(vertices of the triangles) and $\mathcal{S}_{h}=\{s_{k}\}_{k=1}^{E}$ the
\textbf{edges}.

\medskip\noindent
\textbf{Main question (Mesh Generation, Part 1).} Given a collection of points
in $\R^{2}$ (nodes), how can we connect them to form a triangulation?

% ==================================================================
\subsection*{Voronoi diagrams}

\textbf{Idea.} Given a finite set of points in $\R^{2}$, assign to each point a
\emph{region of influence} in such a way that it decomposes $\R^{2}$.

\medskip\noindent
\textbf{Defn (Voronoi region).} Let $X$ be a collection of $N$ points in
$\R^{2}$. For each $x_{i}\in X$ the \textbf{Voronoi region} $V_{i}$ is the set
of points $x\in\R^{2}$ that are at least as close to $x_{i}$ as to any other
point in $X$:
\begin{equation}\label{eq:l13voronoi}
  V_{i}=\bigl\{x\in\R^{2}:\ \|x-x_{i}\|\le\|x-x_{j}\|\ \ \forall\,x_{j}\in X\bigr\} .
\end{equation}
Consider the half-plane
\begin{equation}\label{eq:l13halfplane}
  H_{ij}=\bigl\{x\in\R^{2}:\ \|x-x_{i}\|\le\|x-x_{j}\|\bigr\} .
\end{equation}
The Voronoi region of $x_{i}$ is the convex intersection of the half-planes
$H_{ij}$ for all $j\neq i$; therefore $V_{i}$ is a convex polygonal region
(possibly unbounded).

\medskip\noindent
\textbf{Notes.}
\begin{enumerate}[label=(\arabic*),leftmargin=2.4em,itemsep=2pt]
\item Each point in $\R^{2}$ has at least one nearest point in $X$, so every
  $x\in\R^{2}$ lies in at least one Voronoi region.
\item Two Voronoi regions lie on opposite sides of a perpendicular bisector;
  two Voronoi regions do not share any interior points. If a point belongs to
  more than one region, it must lie on the boundary.
\item Voronoi regions for a given $X$ cover $\R^{2}$.
\item Voronoi regions together with edges and vertices form the
  \textbf{Voronoi diagram} of $X$.
\item For a given $X$, the Voronoi diagram is unique.
\end{enumerate}

% ==================================================================
\subsection*{Delaunay triangulation}

\textbf{Defn (Delaunay partition).} Let $X\subset\R^{2}$ be a set of $N$
points and let $V$ be the Voronoi diagram. The \textbf{Delaunay partition} is
obtained by creating an edge between two points in $X$ \emph{if and only if}
their Voronoi regions share a boundary.

\medskip\noindent
\textbf{Generic situation:} 3 points of $X$ share a common vertex, so a
triangle is formed.

\medskip\noindent
\textbf{Degenerate situation:} 4 or more points share a common Voronoi vertex
--- \ie they are co-circular.

\medskip\noindent
\textbf{Note.} Degeneracies can always be converted to a triangulation (a
square is split by either diagonal).

\medskip\noindent
\textbf{Defn (Delaunay triangulation).} The Delaunay triangulation $\Th$ is the
Delaunay partition together with some strategy to convert degenerate elements
into triangles.

\medskip\noindent
\textbf{Properties.}
\begin{enumerate}[label=(\arabic*),leftmargin=2.4em,itemsep=2pt]
\item The Voronoi diagram and the Delaunay triangulation form a \textbf{duality}.
\item For a given $X\subset\R^{2}$, $V$ is unique, and $\Th$ is unique up to
  degeneracies.
\item $\Th$ contains $\mathcal{O}(N)$ triangles.
\item $\Th$ \textbf{maximises the minimum angle}.
\item A circle circumscribing any triangle in $\Th$ does not contain any other
  point of $X$.
\end{enumerate}

\begin{mechanics}[property (4) is the one you care about]
``Maximises the minimum angle'' is exactly the regularity condition
\eqref{eq:l13regular} in disguise: of all possible triangulations of the same
point set, Delaunay is the one that comes closest to avoiding slivers. That is
why essentially every unstructured mesher starts from a Delaunay triangulation
--- it is not merely convenient, it is optimal for the constant in the
interpolation estimate.

Property (5), the empty-circumcircle condition, is the practical test a code
uses to check or repair a mesh.
\end{mechanics}

% ------------------------------------------------------------------
\subsubsection*{In Python}

\begin{enumerate}[label=\arabic*),leftmargin=2.2em,itemsep=3pt]
\item Create a list of points in $\R^{2}$: \texttt{p[0:N, 0:2]}, where the
  first index is the point index and the second the coordinate $(x,y)$.
\item Create $\Th$ (Delaunay triangulation):
\begin{quote}\ttfamily
import scipy.spatial as sp\\
t = sp.Delaunay(p).simplices
\end{quote}
\item Plot the mesh:
\begin{quote}\ttfamily
import numpy as np\\
import matplotlib.pyplot as plt\\
plt.triplot(p[:,0], p[:,1], t, facecolors=np.zeros(N), edgecolors='k')
\end{quote}
\item Plot a function defined on the nodes:
\begin{quote}\ttfamily
U = ufunc(p)\\
plt.tripcolor(p[:,0], p[:,1], t, U, shading='flat', cmap='inferno')
\end{quote}
\end{enumerate}

\noindent
\textbf{What is \texttt{t}?} $N=$ number of nodes, $M=$ number of triangles;
\texttt{t} is an $M\times3$ array of integers, where \texttt{t[j][i]} gives, for
triangle $j$, the node index $i\in\{0,1,2\}$.

\medskip\noindent
\textbf{Note (area of a triangle).} For vertices $(x_{0},y_{0})$,
$(x_{1},y_{1})$, $(x_{2},y_{2})$, set
$\underline{v}_{1}=(x_{1}-x_{0},\,y_{1}-y_{0},\,0)$ and
$\underline{v}_{2}=(x_{2}-x_{0},\,y_{2}-y_{0},\,0)$; then
\begin{equation}\label{eq:l13area}
  \text{area}=\tfrac12\bigl\|\underline{v}_{1}\times\underline{v}_{2}\bigr\| .
\end{equation}

\medskip\noindent
\textbf{Note.} Delaunay in C: \texttt{http://www.qhull.org}.

% ==================================================================
\subsection*{Mesh generation, Part 2: representing $\Om$ and $\dOm$}

\textbf{Idea.} Describe the domain \emph{implicitly} using a \textbf{level-set
function}.

\medskip\noindent
\textbf{Ex.} $\phi(x)=\|x\|-R$, where $R>0$. Then
\begin{equation}\label{eq:l13levelset}
  \dOm=\bigl\{x\in\R^{2}:\ \phi(x)=0\bigr\}
      =\bigl\{x\in\R^{2}:\ \|x\|=R\bigr\},
  \qquad
  \Om=\bigl\{x\in\R^{2}:\ \phi(x)<0\bigr\},
\end{equation}
a circle of radius $R$ centred at $(0,0)$.

\medskip\noindent
\textbf{Note (convention).}
\begin{equation}\label{eq:l13conv}
  \{\phi=0\}\ \text{is the boundary},\qquad
  \{\phi<0\}\ \text{is the interior of } \Om,\qquad
  \{\phi>0\}\ \text{is the exterior of } \Om .
\end{equation}

\medskip\noindent
\textbf{Note.} For a given $\dOm$ there are many level-set functions; \eg
$\phi(x)=\|x\|^{p}-R^{p}$ for any positive integer $p$ describes the disc of
radius $R$ centred at $(0,0)$.

\medskip\noindent
\textbf{Boolean operations.}
\begin{align}
  \phi&=\min(\phi_{1},\phi_{2}) &&\text{is the \textbf{union} operation: }
    \Om=\Om_{1}\cup\Om_{2},\label{eq:l13union}\\
  \phi&=\max(\phi_{1},\phi_{2}) &&\text{is the \textbf{intersection}
    operation: } \Om=\Om_{1}\cap\Om_{2}.\label{eq:l13inter}
\end{align}

\begin{intuition}[why min is union]
With the sign convention \eqref{eq:l13conv}, ``inside'' means negative. A point
is in $\Om_{1}\cup\Om_{2}$ when it is inside \emph{at least one} of them, \ie
when \emph{at least one} $\phi_{i}$ is negative --- which is exactly
$\min(\phi_{1},\phi_{2})<0$. Intersection needs \emph{both} negative, hence
$\max$. CAD solid modelling is built on these two lines.
\end{intuition}

\begin{dictionary}[mesh generation]
level-set function $\phi$ & implicit CAD geometry; $\phi<0$ is material\\
Voronoi region & the territory nearest a given node\\
Delaunay triangulation & the dual mesh; the ``best'' triangles for a point set\\
$h_{j}/\rho_{j}\le\alpha$ & element aspect-ratio / sliver check\\
\texttt{scipy.spatial.Delaunay} & what a mesher calls underneath\\
\end{dictionary}

\todoq{The scan's \texttt{plt.triplot} call is written with both
\texttt{facecolors} and \texttt{edgecolors} arguments, which belong to
\texttt{tripcolor}, not \texttt{triplot}. Transcribed as written; worth
checking the lecturer's actual code before relying on it.}
